\chapter{Tablas complementarias}\label{ap:complementos}

\section{Coeficientes del corrector poblacional}

La \cref{tab:corrector} trae los coeficientes de la \cref{eq:corrector} para cada variable, ajustados con los
\cifra{datos.sujetos} sujetos. Es el corrector que viaja con \texttt{nova-espacio} en Hugging Face. Ninguna cifra del
\cref{cap:resultados} sale de él: en la validación, cada fold usa un corrector ajustado sin los sujetos que evalúa.
Las dos últimas columnas son el cambio que el corrector espera a los 30 y 60 días en cama, en la unidad de la variable.
El peso baja, el pulso y la presión diastólica suben, y la presión sistólica, la respiración y la temperatura casi no
cambian, que es el mismo orden del desplazamiento de la \cref{sec:desplazamiento}. El agua ingerida refleja la meta del
protocolo más que la fisiología.

\begin{table}[htbp]
\caption[Coeficientes del corrector poblacional]{Coeficientes del corrector poblacional
$\Delta(d) = \beta_0 + \beta_1 \log(1+d) + \beta_2\, d/30$, ajustado con todos los sujetos, y su valor a los 30 y
60 días en cama.}
\label{tab:corrector}
\centering\small
\begin{tabular}{llrrrrr}\toprule
 Variable & Unidad & $\beta_0$ & $\beta_1$ (log(1+d)) & $\beta_2$ (d/30) & $\Delta(30)$ & $\Delta(60)$ \\
\midrule
Peso & kg & 0,528 & −0,524 & 0,069 & −1,20 & −1,49 \\
Presión diastólica & mmHg & −2,539 & 1,042 & 0,148 & 1,19 & 2,04 \\
Pulso & lpm & −4,627 & 1,575 & 0,748 & 1,53 & 3,34 \\
Diuresis & g/día & 63,4 & −75,9 & 117,1 & −80 & −14 \\
Frecuencia respiratoria & rpm & 0,318 & −0,133 & 0,131 & −0,01 & 0,03 \\
Presión sistólica & mmHg & −0,860 & 0,347 & −0,333 & 0,00 & −0,10 \\
Temperatura & °C & 0,044 & −0,016 & 0,040 & 0,03 & 0,06 \\
Agua ingerida & g/día & −144,417 & 8,477 & 13,864 & −101 & −82 \\
\bottomrule\end{tabular}

\end{table}

\section{Pronóstico por variable}

La \cref{tab:pronostico-var} desglosa por variable el pronóstico en cama de la \cref{tab:pronostico}. Salvo en el peso, los
cinco PatchTST quedan a una centésima o menos entre sí. En el peso la media de 7 días le gana a todos, y en la
frecuencia respiratoria gana Chronos-2. En las tres variables de la v4, el pulso y las dos presiones, el mejor es el
PatchTST entrenado desde cero, por milésimas: otra razón para no dar por hecho que el pre-entrenamiento aporta
(\cref{sec:v4}).

\begin{table}[htbp]
\caption[Pronóstico en cama por variable]{MAE del pronóstico a 7 días en cama, relativo a la persistencia, por
variable. En negrita el menor de cada fila, comparado con todos sus decimales. Cero: PatchTST entrenado solo con el
reposo. Sin LS: pre-entrenado sin LifeSnaps. Tierra: sin adaptar. Cabeza y Todo: adaptados desde Tierra.}
\label{tab:pronostico-var}
\centering\small
\begin{tabular}{lrrrrrrr}\toprule
 Variable & Media 7 & Chronos-2 & Cero & Sin LS & Tierra & Cabeza & Todo \\
\midrule
Peso & \textbf{0,880} & 0,899 & 0,916 & 0,900 & 0,909 & 0,908 & 0,892 \\
Presión diastólica & 0,773 & 0,802 & \textbf{0,758} & 0,761 & 0,767 & 0,761 & 0,760 \\
Pulso & 0,776 & 0,790 & \textbf{0,748} & 0,750 & 0,757 & 0,749 & 0,749 \\
Diuresis & 0,758 & 0,777 & 0,744 & 0,746 & 0,749 & 0,743 & \textbf{0,743} \\
Frecuencia respiratoria & 0,870 & \textbf{0,840} & 0,853 & 0,856 & 0,847 & 0,855 & 0,852 \\
Presión sistólica & 0,772 & 0,792 & \textbf{0,753} & 0,755 & 0,763 & 0,756 & 0,756 \\
Temperatura & 0,818 & 0,839 & 0,806 & 0,802 & 0,802 & 0,800 & \textbf{0,799} \\
Agua ingerida & 0,817 & 0,818 & 0,793 & 0,794 & 0,799 & \textbf{0,792} & 0,793 \\
\bottomrule\end{tabular}

\end{table}
